\chapter{Comunicazione Distribuita}

La comunicazione tra agenti avviene tramite un bus condiviso gestito dal
\texttt{CommunicationManager}, un componente centralizzato a livello di modello
che simula trasmissioni locali a raggio limitato. Ogni agente comunica solo con
i vicini entro il proprio \texttt{communication\_radius}: non esiste alcun canale
globale.

\section{Messaggi e Wrapper Timestamped}

Tutti i messaggi scambiati ereditano da \texttt{Message} e contengono un campo
\texttt{step} che registra il passo di simulazione in cui sono stati generati.
Quando un agente riceve dati duplicati da fonti diverse (ad~esempio la posizione
dello stesso Retriever riportata da due Coordinator), viene applicata la regola
\textit{newest wins}: solo l'informazione con \texttt{step} pi\`u recente viene
mantenuta. Questo evita che dati stale sovrascrivano informazioni aggiornate
durante la propagazione multi-hop.

I principali tipi di messaggio sono:

\begin{itemize}
    \item \texttt{MapDataMessage}: Contiene la sottomatrice della mappa esplorata,
          i \texttt{known\_objects}, le posizioni dei Retriever e gli oggetti reclamati.
          Scambiato periodicamente tra tutti i tipi di agente.
    \item \texttt{TaskAssignment}: Compito assegnato dal Coordinator a un Retriever
          specifico, contenente la posizione dell'oggetto.
    \item \texttt{TaskStatusMessage}: Diffuso dai Retriever ai vicini, contiene la
          \texttt{task\_queue} corrente e la posizione dell'agente. Permette ai peer
          di rimuovere task duplicati dalla propria coda.
    \item \texttt{RetrieverEventMessage}: Messaggio peer-to-peer tra Retriever con
          eventi \texttt{object\_spotted} e \texttt{cargo\_dropped}, usato per
          propagare la scoperta di oggetti e il rilascio di carico.
\end{itemize}

\section{Oggetti Reclamati e Protocollo di Takeover}

Il dizionario globale \texttt{claimed\_objects} nel \texttt{CommunicationManager}
associa ogni posizione di oggetto reclamato a una tupla
$(\text{agent\_id}, \text{step})$. La funzione \texttt{try\_claim\_object} implementa
un protocollo atomico con quattro condizioni di successo (in ordine):

\begin{enumerate}
    \item L'oggetto non \`e mai stato reclamato.
    \item L'oggetto \`e gi\`a reclamato dallo stesso agente (idempotenza).
    \item Il claim esistente \`e pi\`u vecchio di \texttt{\_STALE\_CLAIM\_AGE} passi
          (default~50), ed \`e quindi considerato scaduto.
    \item Il richiedente ha energia superiore a $1{,}5\times$ quella del detentore
          corrente (preemption energetica).
\end{enumerate}

Se nessuna condizione \`e soddisfatta, il claim viene rifiutato. Il meccanismo di
stale claim takeover impedisce che un agente bloccato o fuori energia monopolizzi
un oggetto indefinitamente, mentre la preemption energetica favorisce l'agente
con maggiore autonomia residua.

\section{Contatore dei Messaggi}

Il \texttt{CommunicationManager} mantiene un contatore \texttt{messages\_sent}
incrementato per ogni destinatario effettivo di un messaggio. Se un agente invia
un messaggio e questo raggiunge $n$~agenti nel raggio di comunicazione, il contatore
viene incrementato di~$n$. Questo valore viene esposto nelle metriche di simulazione
e nel benchmark per misurare il costo comunicativo dello sciame.

\section{Metriche di Efficienza Energetica}

Le metriche del sistema includono il rapporto di efficienza energetica:

$$\eta = \frac{n_{\text{delivered}}}{E_{\text{total\_consumed}}}$$

dove $n_{\text{delivered}}$ \`e il numero di oggetti consegnati e $E_{\text{total\_consumed}}$
\`e l'energia totale consumata da tutti gli agenti. Un valore pi\`u alto indica un uso
pi\`u efficiente dell'energia per unit\`a di lavoro utile.
